% Ch13 WS3 -- Le Chatelier + FRQ. Block 5.
\documentclass{apchem-worksheet}

\apcunitword{Chapter}
\apcunit{13}
\apcunittitle{Chemical Equilibrium}
\apcdoctitle{Worksheet 3 \textbullet\ Le Ch\^{a}telier's Principle \& FRQ}
\apcsubtitle{Zumdahl \S13.7 \textbullet\ only temperature changes the value of $K$}

\begin{document}
\wsheader[For every prediction, state the direction of shift AND whether the
value of $K$ changes. Those are two separate questions and both are
examinable.]

\problem[]{For \ce{2 SO2(g) + O2(g) <=> 2 SO3(g)},
$\Delta H = \SI{-198}{\kilo\joule}$, predict the shift and state whether $K$
changes:}
\begin{parts}
  \item Add \ce{SO2}: \answerblank[5cm]{right; $K$ unchanged}
  \item Remove \ce{SO3}: \answerblank[5cm]{right; $K$ unchanged}
  \item Decrease the volume: \answerblank[7.4cm]{right (3 mol $\to$ 2 mol);
        $K$ unchanged}
  \item Increase the temperature:
        \answerblank[5cm]{left; $K$ \emph{decreases}}
  \item Add a catalyst: \answerblank[5cm]{no shift; $K$ unchanged}
  \item Add helium at constant volume:
        \answerblank[5cm]{no shift; $K$ unchanged}
\end{parts}

\problem[]{Explain each of the three "no shift" cases above:}
\begin{parts}
  \item Why does a catalyst cause no shift?
        \answerlines[2]{It lowers the activation energy for both directions
        equally, so equilibrium is reached faster but its position is
        unchanged.}
  \item Why does adding an inert gas at constant volume cause no shift?
        \answerlines[2]{The partial pressures of the reacting gases are
        unchanged, so $Q$ still equals $K$ --- only the total pressure
        rises.}
  \item Why does adding more of a pure solid cause no shift?
        \answerlines[2]{Solids do not appear in $K$; their concentration is
        fixed by density, so adding more changes no term in the
        expression.}
\end{parts}

\problem[]{Only one kind of change alters the numerical value of $K$.}
\begin{parts}
  \item Name it. \answerblank[3cm]{temperature}
  \item Explain why the others do not. \answerlines[3]{Any other disturbance
        changes concentrations away from equilibrium, and the system shifts
        precisely far enough to restore the same value of $K$. Temperature
        changes the underlying thermodynamics, so the constant itself takes
        a new value.}
  \item For an endothermic reaction, does raising the temperature increase
        or decrease $K$? Explain using heat as a reactant.
        \answerlines[2]{Increase. Writing heat as a reactant, adding heat
        pushes the system to the right, producing more products at
        equilibrium --- a larger $K$.}
\end{parts}

\problem[]{Zumdahl's arsenic roasting reaction:
\ce{As4O6(s) + 6 C(s) <=> As4(g) + 6 CO(g)}. Predict the effect of each
change:}
\begin{parts}
  \item Adding \ce{CO}: \answerblank[3cm]{shifts left}
  \item Adding more carbon: \answerblank[4.2cm]{no shift (pure solid)}
  \item Removing \ce{As4} gas as it forms:
        \answerblank[3cm]{shifts right}
  \item Increasing the container volume:
        \answerblank[5.9cm]{shifts right (more gas moles)}
\end{parts}

\problem[]{\textbf{FRQ (10 points).} The Haber process synthesizes ammonia:
\[ \ce{N2(g) + 3 H2(g) <=> 2 NH3(g)} \qquad \Delta H = \SI{-92}{\kilo\joule} \]}
\begin{parts}
  \item Write the $K_c$ expression, and determine $\Delta n$ and the
        relationship between $K_p$ and $K_c$.
        \answerlines[2]{$K_c = [\ce{NH3}]^2/([\ce{N2}][\ce{H2}]^3)$;
        $\Delta n = 2-4 = -2$, so $K_p = K_c(RT)^{-2}$.}
  \item A reaction vessel at equilibrium contains
        $[\ce{N2}] = \SI{0.50}{\Molar}$, $[\ce{H2}] = \SI{1.0}{\Molar}$,
        and $[\ce{NH3}] = \SI{0.20}{\Molar}$. Calculate $K_c$.
        \workspace[2cm]
        \keyonly{\begin{apcanswer}$K_c = (0.20)^2/[(0.50)(1.0)^3] =
        0.040/0.50 = 0.080$\end{apcanswer}}
  \item Some \ce{NH3} is then removed. State the direction of shift, and
        state whether $K_c$ changes. \answerlines[2]{The system shifts right
        to replace the ammonia. $K_c$ is unchanged --- only temperature
        alters it.}
  \item Predict and explain the effect on the equilibrium yield of \ce{NH3}
        of (i) increasing the pressure and (ii) increasing the temperature.
        \answerlines[4]{(i) Higher pressure shifts the system toward the
        side with fewer gas moles --- 2 on the right versus 4 on the left ---
        so the yield increases. (ii) The forward reaction is exothermic, so
        adding heat shifts the system left; the yield decreases and $K$
        itself becomes smaller.}
  \item Industrial plants nevertheless operate at roughly
        \SI{450}{\celsius}. Explain this apparent contradiction.
        \answerlines[3]{At low temperature the equilibrium position is
        favourable but the rate is impractically slow. Running hot sacrifices
        some equilibrium yield to reach equilibrium quickly; high pressure
        and a catalyst recover much of the loss. It is a deliberate
        kinetics-versus-thermodynamics compromise.}
\end{parts}

\begin{apcnote}[Rubric --- score yourself, 10 points]
\textbf{(a) 2 pts} 1 $K_c$ expression, 1 $\Delta n$ with the $K_p$ relation
\textbullet\ \textbf{(b) 1 pt} 0.080 \textbullet\ \textbf{(c) 2 pts}
1 direction, 1 explicitly states $K$ unchanged \textbullet\
\textbf{(d) 3 pts} 1 pressure with the mole-count reason, 1 temperature
direction, 1 states that $K$ itself decreases \textbullet\
\textbf{(e) 2 pts} 1 identifies the rate problem, 1 names the trade-off
against yield.
\end{apcnote}

\begin{spiralstrip}{Chapter 13 \textbullet\ blocks 1--4}
  \item If $Q > K$, the shift is: \answerblank[2cm]{left}
  \item $K$ expression for \ce{C(s) + CO2(g) <=> 2 CO(g)}:
        \answerblank[4cm]{$[\ce{CO}]^2/[\ce{CO2}]$}
  \item The 5\% rule checks: \answerblank[6.0cm]{validity of the
        approximation}
  \item Does a large $K$ mean fast? \answerblank[1.6cm]{no}
  \item $\Delta n$ for \ce{2 NO2 <=> N2O4}: \answerblank[1.6cm]{$-1$}
\end{spiralstrip}

\end{document}
